% ===== CHAPTER 1: INTRODUCTION =====

\chapter{Introduction}

\section{Background and Context}

Planning a trip typically requires a traveller to move between many
disconnected services: one source for destination ideas, another for
attractions, another for accommodation, a weather service, a map, and
separate sources for health and safety advisories. When the trip is planned
with friends or family, the group usually adds a messaging application and a
shared spreadsheet to compare options and reach a decision, while tickets and
booking confirmations remain scattered across e-mail inboxes and files. The
result is a planning process that is time-consuming, fragmented and difficult
to personalise.

TravelSync is a web application that brings these activities into a single
platform for travel within India. It serves two kinds of travellers. A
traveller who already knows the destination can select it directly and move
straight to exploring and preparing for it. A traveller who is undecided can
instead describe what the trip should offer, namely travel interests, the
importance of each interest, traveller type, trip duration, total budget,
starting location and the maximum distance they are willing to travel, and
receive a ranked list of destinations that are both similar to these
preferences and practically feasible.

Once a destination is chosen it becomes a \emph{trip}, and every other
module of the platform is organised around that trip: the five top places at
the destination, a smart packing list that adapts to the weather forecast,
accommodation recommendations near the places the traveller intends to visit,
destination-specific health and safety information, an interactive map with
driving distances, and a store for bookings and tickets. A Collaborative
Travel Hub lets a group create a shared space, chat in real time, suggest
destinations, vote, and lock in a final decision that can be converted
directly into a trip. A floating TravelSync assistant answers travel questions
from the platform's own data throughout the application.

\subsection{Problem Statement}

Existing travel-planning workflows are scattered across multiple platforms
and are rarely personalised enough to be genuinely useful. Specifically:

\begin{itemize}
  \item Travellers who are unsure where to go usually have implicit
    constraints, such as budget, duration and acceptable travel distance,
    but no single tool converts these constraints together with their
    interests into a ranked and comparable set of destination options.
  \item General-purpose recommendation lists rarely consider a traveller's
    interests, traveller type, trip duration, budget and starting location
    at the same time.
  \item Information needed once a destination is chosen, such as places to
    visit, accommodation, weather and health advisories, is distributed
    across different services instead of being available in one place.
  \item Packing is handled with generic checklists that do not account
    together for the destination's climate and terrain, the trip length and
    the forecast for the actual travel dates.
  \item Groups travelling together lack a structured way to propose
    destinations, discuss them and reach a collective decision, and fall
    back on separate communication tools.
  \item Bookings, tickets and records of past trips are spread across
    e-mails, files and applications rather than being attached to the trip
    or group they belong to.
\end{itemize}

TravelSync addresses these problems by combining a content-based
recommendation engine with deterministic feasibility constraints,
consolidated destination information, trip-preparation tools, real-time group
collaboration with voting, and booking and travel-history management within
one secure, account-based platform.

\subsection{Existing Systems and Proposed System}

Travellers today rely on several categories of existing systems, each of
which solves part of the problem. Online travel agencies and hotel
aggregators are optimised for booking transport and accommodation once a
destination is already known; they do not help an undecided traveller choose
where to go. Map and review platforms describe places in detail but treat
each destination in isolation. Weather services, public-health advisories
and packing checklists are separate tools, and group decisions are typically
made in general-purpose messaging applications in which suggestions, votes
and the final decision are lost in the conversation history.

The proposed system differs from these in three ways. First, it supports the
\emph{undecided} traveller with an explainable, content-based recommender
\cite{lops2011content} that ranks destinations by similarity to the
traveller's interests and then applies duration, budget and distance
constraints. Second, it makes the trip the central object: every supporting
module (places, packing, stays, health, map and bookings) reads its context
from the saved trip, so the traveller never re-enters the destination, dates
or budget. Third, it treats group planning as a first-class workflow, with
structured destination suggestions, one vote per member, live updates and an
owner-confirmed decision that becomes a shared trip.

\section{Project Objectives}
\label{sec:objectives}

\begin{enumerate}
  \item \textbf{Objective 1:} Build an intelligent trip planner that lets a
    user either select a known destination or obtain a ranked list of up to
    ten recommended destinations from a dataset of 250 Indian destinations,
    using weighted nearest-neighbour similarity over twelve interest features
    and four traveller types, combined with duration, budget and
    great-circle distance constraints.
  \item \textbf{Objective 2:} Consolidate the information and preparation a
    chosen destination needs into a single trip workspace: the top five
    places, a weather-aware packing list, scored accommodation
    recommendations, destination-specific health and safety information and
    an interactive map with driving distances.
  \item \textbf{Objective 3:} Provide a Collaborative Travel Hub in which
    groups of travellers can join through an invite code, chat in real time,
    suggest destinations, vote, and confirm a final destination that is
    converted into a trip for the group.
  \item \textbf{Objective 4:} Support trip management end to end, through
    bookings and tickets with private file attachments, a record of upcoming
    and previous trips, in-app notifications and a conversational travel
    assistant, while protecting every user's data with secure authentication
    and per-user authorisation.
\end{enumerate}

These objectives were framed to be Specific, Measurable, Achievable, Relevant
and Time-bound (SMART). Each objective maps to concrete modules and API
endpoints described in Chapter~\ref{chap:methodology}, each is verified by
automated tests reported in Chapter~\ref{chap:results}, and all four were
delivered within the semester's development window (August to November
2026).

\subsection{Scope of the Project}

The scope of TravelSync is domestic travel within India. The destination
dataset covers 250 destinations, the places dataset contains the top five
places for each of them (1{,}250 places), and 258 starting locations are
available for distance calculations. The following capabilities are in
scope: account management; known-destination selection and preference-based
recommendation; trip creation, editing and deletion; the top places of each
destination; packing lists; accommodation recommendations; health and safety
information; group collaboration; bookings with file attachments; an
interactive map; notifications; and the travel assistant.

The following are deliberately out of scope. TravelSync does not book or pay
for transport or accommodation; bookings are records that the user enters.
It does not display live hotel prices or availability, because the
accommodation data source does not provide them. Its health module provides
general travel awareness only and does not offer medical advice or a
``safe to travel'' score. There is no administrator interface; the curated
datasets are maintained by the development team as versioned files in the
repository.

Development followed the time-to-value order set out in the project
proposal. The core discovery workflow (known-destination selection and
recommendation) was delivered first, followed by the trip workspace and its
supporting modules (top places, packing, stays and health), then the
Collaborative Travel Hub with voting, and finally bookings, the interactive
map and the travel assistant.

\subsection{Software Requirements Specification}

The requirements below follow the structure recommended by IEEE Std 830
\cite{ieee830}. Table~\ref{tab:functional-requirements} lists the functional
requirements and the module that implements each one.

\begin{table}[htbp]
\centering
\begin{tabular}{|l|L{0.56\textwidth}|l|}
  \hline
  \textbf{ID} & \textbf{Functional requirement} & \textbf{Module} \\
  \hline
  FR-01 & Register, log in, log out, change and reset a password & Auth \\
  \hline
  FR-02 & Search the destination dataset and select a known destination
    & Planner \\
  \hline
  FR-03 & Recommend ranked destinations from interests, importance,
    traveller type, duration, budget, start and maximum distance & Planner \\
  \hline
  FR-04 & Save, list, edit (dates, budget, duration) and delete trips
    & Trips \\
  \hline
  FR-05 & Show the top five places of a destination & Places \\
  \hline
  FR-06 & Generate, edit and tick off a packing list; refine it with a
    forecast once dates are set & Packing \\
  \hline
  FR-07 & Recommend, explain, save and remove accommodation for a trip
    & Stays \\
  \hline
  FR-08 & Show current health alerts and historical outbreak patterns for
    the destination's state & Health \\
  \hline
  FR-09 & Create and join groups, chat in real time, suggest destinations,
    vote and confirm a decision & Groups \\
  \hline
  FR-10 & Record personal and group bookings and attach ticket files
    & Bookings \\
  \hline
  FR-11 & Show a trip's places and saved stays on a map with driving
    distances & Map \\
  \hline
  FR-12 & Answer travel questions through a conversational assistant
    & Assistant \\
  \hline
  FR-13 & Notify users of group events and let them mark notifications read
    & Notifications \\
  \hline
\end{tabular}
\caption{Functional requirements}
\label{tab:functional-requirements}
\end{table}

The non-functional requirements are as follows:

\begin{itemize}
  \item \textbf{Security:} passwords are stored only as bcrypt hashes;
    sessions use a signed JWT in an \texttt{httpOnly} cookie; every private
    endpoint verifies that the requested trip, group, booking or file
    belongs to the caller; and state-changing requests from other origins
    are rejected.
  \item \textbf{Reliability:} failures of external providers (accommodation,
    routing, weather and the optional language model) degrade to a clear
    message or a fallback rather than breaking the page, and a database
    outage returns HTTP 503 instead of logging the user out.
  \item \textbf{Performance:} interactive API requests other than
    deliberately slow password hashing should complete well within the
    100\,ms range on a single server process (measured in
    Section~\ref{sec:performance}).
  \item \textbf{Usability and accessibility:} validation errors are shown
    next to the field concerned, live regions announce results to screen
    readers, and the interface adapts to desktop and mobile widths.
  \item \textbf{Maintainability:} the backend is split into routers,
    services and models, every rule-based engine is a pure, unit-tested
    function, and configuration is supplied only through environment
    variables.
  \item \textbf{Transparency:} recommendation and accommodation scores are
    explained to the user, estimated prices are labelled as estimates, and
    health information is labelled as historical, not predictive.
\end{itemize}

The software required to build and run the system is Python 3.10 or later
with the packages listed in \texttt{backend/requirements.txt}; Node.js
22.12 or later for the frontend build and test tool chain; MongoDB (a local
server or MongoDB Atlas); and a current web browser. Optional components are
a Geoapify API key for accommodation and driving distances, an SMTP server
for password recovery, and an OpenAI-compatible chat API key for the
assistant's fallback answers. No special hardware is required: the
measurements in Chapter~\ref{chap:results} were taken on a standard laptop,
and the client runs in any modern desktop or mobile browser.
